\chapter{Kết luận và Hướng phát triển}
\label{chp:conclusion}

\section{Kết luận}

\subsection{Đạt được}
Đồ án đã nghiên cứu và triển khai thành công một hệ thống vi mạch nhúng SoC (System-on-Chip) hoàn chỉnh dựa trên phương pháp đồng thiết kế Phần cứng -- Phần mềm (Hardware-Software Co-design). Hệ thống được hiện thực hóa trên nền tảng FPGA Xilinx Kria KV260 (Zynq UltraScale+), hướng tới mục tiêu gia tốc các thuật toán Trí tuệ nhân tạo tại biên (Edge AI).

Thông qua quá trình thực hiện, nhóm sinh viên đã đạt được những kết quả cốt lõi sau:
\begin{itemize}
    \item \textbf{Xây dựng thành công kiến trúc xử lý Đa nhân (Dual-core):} Kết hợp linh hoạt giữa sức mạnh quản lý cấp cao của lõi ARM Cortex-A53 (chạy hệ điều hành Linux + thư viện PYNQ) và khả năng xử lý thời gian thực, tiêu thụ năng lượng thấp của vi điều khiển mã nguồn mở RISC-V CV32E40P (chạy bare-metal). Cả hai lõi cùng truy xuất bộ nhớ chung qua kiến trúc dual-port BRAM mà không gây tranh chấp.
    \item \textbf{Tối ưu hóa kiến trúc bộ nhớ Harvard trên FPGA:} Quy hoạch thành công hệ thống bộ nhớ Dual-Port BRAM với 64~KB I-BRAM (Bộ nhớ lệnh) và 256~KB D-BRAM (Bộ nhớ dữ liệu). Việc thiết lập 4 khối AXI BRAM Controller độc lập đã loại bỏ triệt để hiện tượng nghẽn cổ chai (bottleneck) khi ARM và RISC-V cùng truy xuất bộ nhớ. Tổng cộng 99 RAMB36E2 được sử dụng (68.75\% ngân sách BRAM của KV260) -- chấp nhận được, vẫn còn dư cho các mở rộng tương lai.
    \item \textbf{Khắc phục độ trễ dữ liệu bằng AXI HPC + DMA:} Triển khai thành công bộ điều khiển truy cập bộ nhớ trực tiếp (AXI DMA) kết nối qua cổng \texttt{S\_AXI\_HPC0\_FPD} (cache-coherent qua CCI-400). Điều này giúp tối đa hóa băng thông truyền tải ảnh/trọng số từ DDR4 vào pipeline streaming và tự động đồng bộ hóa bộ đệm (Cache Coherency) mà không cần can thiệp phức tạp bằng phần mềm.
    \item \textbf{Lượng tử hóa và chuẩn hóa mô hình CNN:} Thực hiện thành công kỹ thuật Lượng tử hóa sau huấn luyện (Post-Training Quantization -- PTQ), ép kiểu Tiny-VGG từ dấu phẩy động 32-bit (FP32) xuống số nguyên 8-bit (INT8) qua TFLite Interpreter. Trên tập \textit{Cats vs Dogs}, mô hình đạt độ chính xác 79.38\% (FP32) / $\sim$78\% (INT8) -- mức rớt $\leq$~1.5\% nằm trong ngưỡng an toàn theo nghiên cứu của Jacob~et~al.~\cite{jacob2018quantization}, chứng minh PTQ là hợp lý cho cấu hình mạng nhỏ này.
    \item \textbf{Bộ gia tốc Conv-CNN v2.0 hoạt động đúng chuẩn TFLite:} Thiết kế từ đầu một datapath INT8 đầy đủ gồm line buffer 2-bank parity-toggled, window 3$\times$3, 9 PE MAC pipelined (DSP48E2-mappable), bộ requantize 3-stage (bias add + Q31 multiply + arithmetic shift + zero-point + saturate + ReLU). Tài nguyên thực tế: 9 DSP (0.72\%), 22{,}491 LUT (19.20\%), 28{,}715 FF (12.26\%) -- chứng tỏ thiết kế tuần tự rất tiết kiệm tài nguyên, để dành dư địa cho các mở rộng V3.0.
    \item \textbf{Tổng hợp Vivado đạt timing và năng lượng tốt:} Worst Negative Slack (WNS) = +2.118~ns $\Rightarrow$ Fmax thực tế $\sim$127~MHz (đề tài chốt 100~MHz cho đồng bộ với clock của PS). Tổng công suất ước tính 2.945~W, trong đó toàn bộ logic người dùng PL (RISC-V + Conv-CNN + AXI) chỉ tốn $\sim$0.225~W -- riêng bộ gia tốc \texttt{conv\_cnn\_0} chỉ tiêu thụ 23~mW, minh chứng cho hiệu quả năng lượng của giải pháp loosely-coupled trên FPGA.
\end{itemize}

Nhìn chung, đồ án đã chứng minh tính khả thi và hiệu quả của việc tự thiết kế một bộ gia tốc AI chuyên dụng (Custom AI Accelerator) giao tiếp với softcore RISC-V trên nền tảng FPGA, thay vì phụ thuộc hoàn toàn vào các vi xử lý thương mại có sẵn (Vitis AI DPU, Edge TPU). So với các framework liên quan (Vitis AI DPU, FINN, hls4ml, CFU Playground), đề tài có vị trí riêng ở phân khúc \textit{open-source, loosely-coupled, INT8, có RISC-V orchestrator} -- chưa được lấp đầy trong giai đoạn nghiên cứu hiện tại.

\subsection{Hạn chế của hệ thống}
\begin{itemize}
    \item \textbf{Phạm vi op còn hạn chế:} Conv-CNN v2.0 hiện chỉ hỗ trợ Conv 3$\times$3 stride~1 với padding \texttt{valid} và activation ReLU/None. Chưa hỗ trợ Conv 1$\times$1, depthwise separable, pooling, batch normalization, padding \texttt{same}, hay residual add. Điều này giới hạn các mô hình deploy được hiện tại ở mức Tiny-VGG; các kiến trúc phổ biến hơn như VGG11, ResNet18, MobileNet chưa thể chạy trực tiếp trên board.
    \item \textbf{Per-tensor quantization:} TFLite hỗ trợ per-axis quantization (mỗi output channel có scale/zero\_point riêng), giúp giảm rớt accuracy thêm 1--2\%. Đề tài hiện dùng per-tensor (lấy \texttt{scales[0]}) để đơn giản hóa RTL requantize. Đây là trade-off chấp nhận được cho prototype, nhưng nên nâng cấp ở phiên bản sản xuất.
    \item \textbf{Throughput tuần tự:} Kiến trúc 9-PE chạy channel-interleaved + serial \texttt{cnt\_cout} cho latency $\sim$520~ms/frame ($\sim$1.9~FPS) ở 100~MHz -- chưa đạt mức real-time (15--30~FPS). Bottleneck chính là Layer~3 (chiếm 67\% latency do có 32 input channel $\times$ 64 output channel).
    
    \item \textbf{Mới hỗ trợ phân loại nhị phân:} Đề tài giới hạn ở task classification 2 lớp (cats\_dogs). Tập INRIA Person ở định dạng PASCAL-VOC chưa được khai thác do yêu cầu pipeline detection (parse XML + bounding box) vượt scope hiện tại.
\end{itemize}

\newpage
\section{Hướng phát triển}
Mặc dù hệ thống đã hoàn thành các mục tiêu cốt lõi, đồ án vẫn còn nhiều không gian để tối ưu hóa và mở rộng. Hướng phát triển được phân chia theo 3 tầng độ ưu tiên (Tier~A: tăng hiệu năng, Tier~B: mở rộng task, Tier~C: nâng tầm hệ điều hành / UI).

\subsection*{Tier~A -- Nâng cấp hiệu năng inference (cho Tiny-VGG hiện tại)}
\begin{itemize}
    \item \textbf{Conv-CNN v3.0 -- Parallel cout:} Mở rộng từ 9~PE đơn lẻ lên $9\times N$ PE chạy song song trên trục \texttt{cnt\_cout}. Layer~3 (bottleneck) sẽ giảm latency theo tỉ lệ $1/N$. Để đạt 30~FPS cần $N \geq 16$ (tương đương 144 DSP, chỉ 11.5\% ngân sách 1{,}248 DSP của KV260). RTL change tập trung ở \texttt{controller\_fsm.sv} (counter mới) và \texttt{conv\_core.sv} (mảng PE + adder tree mở rộng).
    \item \textbf{Tối ưu luồng dữ liệu bằng Ping-Pong + Tiling:} Mở rộng logic điều khiển DMA và RISC-V để áp dụng kỹ thuật chia gạch (Tiling) và đệm bóng bàn (Ping-Pong Buffering) trong D-BRAM. Điều này sẽ giúp che giấu hoàn toàn độ trễ truyền dữ liệu (Latency Hiding), cho phép khối CNN tính toán liên tục không có thời gian chết.
    \item \textbf{Per-axis quantization:} Cập nhật \texttt{requantize.sv} để chứa mảng $M_{q31}$ và \texttt{shift} per-output-channel thay vì một giá trị duy nhất. Bộ nhớ tăng thêm $\sim$1~KB BRAM nhưng có thể buy thêm 1--2\% accuracy.
\end{itemize}

\subsection*{Tier~B -- Mở rộng task và op support}
\begin{itemize}
    \item \textbf{MaxPool 2$\times$2 + PAD\_SAME (Phase~B):} Module \texttt{max\_pool.sv} đã có sẵn, chỉ cần wire vào pipeline + verify simulation. PAD\_SAME đòi hỏi mở rộng \texttt{controller\_fsm.sv} để zero-pad ở edge. Đây là điều kiện tiên quyết để chạy được VGG11 trên FPGA.
    \item \textbf{Conv 1$\times$1 + Skip-add + BatchNorm fold (cho ResNet18):} Conv 1$\times$1 cần MUX trong \texttt{window\_3x3.sv}; skip-add đòi hỏi 1 port AXIS thứ hai cùng adder ở \texttt{conv\_core.sv}; BN có thể fold vào weight + bias offline (TFLite tự fold khi PTQ).
    \item \textbf{Phát hiện đối tượng (Object Detection -- Phase~C):} Mở rộng pipeline để chạy YOLO-tiny / SSD-Lite. Cần thêm: anchor boxes, multi-scale feature pyramid, non-max suppression (NMS) trên ARM, leaky-ReLU, upsample 2$\times$ nearest. Tập INRIA Person ở định dạng PASCAL-VOC sẽ được tận dụng đầy đủ ở giai đoạn này.
\end{itemize}

\subsection*{Tier~C -- System-level + UX}
\begin{itemize}
    \item \textbf{Mở rộng tập lệnh RISC-V (Custom Instructions):} Tận dụng đặc tính mã nguồn mở của RISC-V để can thiệp vào bộ giải mã lệnh (Instruction Decoder), tạo ra các tập lệnh tùy chỉnh dành riêng cho việc cấu hình và kích hoạt khối gia tốc CNN thay vì dùng giao tiếp Memory-Mapped I/O thông thường. Điều này có thể giảm overhead mỗi layer khoảng 5--10 chu kỳ.
    \item \textbf{Tích hợp Hệ điều hành thời gian thực (RTOS):} Nâng cấp phần mềm chạy trên lõi RISC-V từ dạng bare-metal (vòng lặp vô tận) lên một hệ điều hành thời gian thực nhỏ gọn như FreeRTOS. Điều này sẽ giúp hệ thống có khả năng quản lý nhiều tác vụ đồng thời (nhận dữ liệu cảm biến, điều khiển động cơ, chạy AI) với độ trễ phản hồi xác định.
    \item \textbf{Camera real-time + Web UI:} Tận dụng hệ điều hành Ubuntu chạy trên lõi ARM để phát triển một máy chủ Web cục bộ. Máy chủ này tiếp nhận luồng video trực tiếp từ USB Camera (qua V4L2) hoặc MIPI camera native của Kria, gọi module gia tốc phần cứng dưới PL để suy diễn, và hiển thị kết quả Bounding Box trực quan lên trình duyệt theo thời gian thực. Khi V3.0 đạt $\geq$15~FPS, demo này sẽ mang tính trình diễn cao cho ứng dụng smart camera.
\end{itemize}
